\section{项目简介}

本项目是一个 \textbf{STM32F407} 系列微控制器的裸机（Bare-metal）工程，
使用 \textbf{CMake} 组织构建，配合 \textbf{arm-none-eabi-gcc} 交叉工具链，
采用 ST 官方较老的 \textbf{标准外设库（StdPeriph Driver）} 而非 HAL 库。

\subsection{硬件目标}

\begin{itemize}
  \item 芯片：STM32F407VET6
  \item 内核：ARM Cortex-M4（带 FPU），主频最高 168\,MHz
  \item Flash：512\,KB；RAM：192\,KB（含 64\,KB CCM）
  \item 当前演示功能：GPIO 点亮/翻转 PD12 上的 LED（跑马灯）
\end{itemize}

\subsection{软件栈}

\begin{table}[h]
  \centering
  \begin{tabular}{ll}
    \toprule
    组件 & 说明 \\
    \midrule
    构建系统 & CMake（$\geq 3.24$），Unix Makefiles / Ninja / MinGW \\
    交叉编译器 & arm-none-eabi-gcc（本项目用 GNU 16.2.0） \\
    库 & ST 标准外设库 StdPeriph + ARM CMSIS \\
    链接脚本 & linker/STM32F407VET6.ld \\
    启动文件 & startup/startup\_stm32f40\_41xxx.s \\
    文档 & 本文档（LaTeX，ctex + xelatex） \\
    \bottomrule
  \end{tabular}
\end{table}

\subsection{工程特点}

\begin{itemize}
  \item \textbf{纯 C + 汇编}，不依赖任何 IDE 工程文件（无 Keil/IAR 工程）。
  \item 使用 \lstinline|-nostdlib -nostartfiles|，完全不链接标准 C 库，
        启动流程由自己的 \texttt{startup} 文件完成。
  \item 工具链文件（\texttt{cmake/arm-none-eabi-gcc.cmake}）与构建脚本解耦，
        跨平台（macOS / Linux / Windows）复用同一份 CMake。
  \item 应用代码（\texttt{User/}）与 ST 设备支持代码（\texttt{Core/}）、
        厂商库（\texttt{Driver/}）分层清晰。
\end{itemize}

\subsection{为什么用 CMake 而不是 IDE}

传统 STM32 开发通常依赖 Keil MDK / IAR / STM32CubeIDE 等 IDE，
优点是开箱即用，缺点是工程文件与 IDE 绑定、难以做持续集成（CI）、
难以在命令行自动化构建。使用 CMake 的好处是：

\begin{enumerate}
  \item 一份 \texttt{CMakeLists.txt} 跨三平台复用，无需手工维护多个工程文件；
  \item 构建过程可脚本化，方便接入 CI / 自动化测试；
  \item 依赖关系显式、可追溯，新增/删除源文件只需改一处；
  \item 不绑定具体编辑器，VSCode / CLion / Vim 均可配合使用。
\end{enumerate}
